% ======================================================================
%  SIMC 2026 --- PRELIMINARY WORKING REPORT  \textperiodcentered{}  Questions 1, 2 & 3
%  Endeavour Track  \textperiodcentered{}  INTERNAL DRAFT --- full proofs, all details.
%
%  Source of all mathematical content:
%     notes/Note 26 May 2026.pdf      (Q1, parts a--b)
%     notes/Note 26 May 2026 (2).pdf  (Q1, parts a--e, g)
%     notes/task 2.pdf                (Q2, parts a--e)
%     src/task1/task1.ipynb           (Q1, parts f, h --- numerical validation)
%     src/task_3.ipynb                (Q3, parts a--g --- data analysis)
%
%  Every deviation from those notes (mistake fixed, gap filled, extra
%  added, elaboration) is logged in  changelog.md  in this folder.
%
%  Team:  K. Vishnu Kumar \textperiodcentered{} Adithya Kesan Jayakanth \textperiodcentered{} Adithya Anand
%  Date:  27 May 2026  (Q3 added; Q1/Q2 unchanged from the 26 May draft)
% ======================================================================

\documentclass[12pt,a4paper]{article}

\usepackage[a4paper, margin=2.0cm, headheight=14pt]{geometry}
\usepackage{mathptmx}                 % Times for text + math
\usepackage[T1]{fontenc}
\usepackage{amsmath,amssymb,amsthm,bm}
\usepackage{enumitem}
\usepackage{xcolor}
\usepackage{fancyhdr}
\usepackage{graphicx}
\usepackage{booktabs}
\usepackage[font=small,skip=3pt]{caption}
\usepackage[hidelinks]{hyperref}

% Q3 figures: space-free copies live in doc/FINAL/figs/; originals in data/output/TASK 3/.
\graphicspath{{../figs/}}

\linespread{1.10}
\setlength{\parskip}{5pt}
\setlength{\parindent}{0pt}

% --- header / footer --------------------------------------------------
\pagestyle{fancy}
\fancyhf{}
\fancyhead[L]{\small SIMC 2026 \textperiodcentered{} Preliminary Working Report (Q1, Q2 \& Q3)}
\fancyhead[R]{\small \thepage}
\fancyfoot[C]{\small Vishnu \textperiodcentered\ Adithya K.J. \textperiodcentered\ Adithya A. \quad\textbf{INTERNAL DRAFT --- NOT FOR SUBMISSION}}
\renewcommand{\headrulewidth}{0.4pt}

% --- review flags (mirrored in extras.md) -----------------------------
\newcommand{\miss}[1]{{\textbf{\textcolor{red}{[MISS: #1]}}}}
\newcommand{\fix}[1]{{\textbf{\textcolor{orange!85!black}{[CORRECTED: #1]}}}}
\newcommand{\extra}[1]{{\textbf{\textcolor{blue!80!black}{[EXTRA: #1]}}}}
\newcommand{\verified}[1]{{\textbf{\textcolor{green!45!black}{[VERIFIED: #1]}}}}

% --- proof environments -----------------------------------------------
\theoremstyle{definition}
\newtheorem*{claim}{Claim}
\newtheorem*{prop}{Proposition}
\DeclareMathOperator{\tr}{tr}

% --- notation shortcuts -----------------------------------------------
\newcommand{\R}{\mathbb{R}}
\newcommand{\Mx}{\mathbf{M}}
\newcommand{\Xb}{\mathbf{X}}
\newcommand{\vv}{\mathbf{v}}
\newcommand{\uu}{\mathbf{u}}
\newcommand{\xx}{\mathbf{x}}
\newcommand{\yy}{\mathbf{y}}
\newcommand{\pp}{\mathbf{p}}
\newcommand{\rr}{\mathbf{r}}
\newcommand{\zero}{\mathbf{0}}
\newcommand{\norm}[1]{\lVert #1 \rVert}
\newcommand{\T}{^{\!\top}}
% Q3 notation
\newcommand{\Xt}{\widetilde{\mathbf X}}
\newcommand{\xt}{\widetilde{\mathbf x}}
\newcommand{\Nt}{\widetilde{N}}
\newcommand{\Cc}{\mathbf{C}}
\newcommand{\Gg}{\mathbf{G}}
\newcommand{\nuv}{\boldsymbol{\nu}}

\title{\vspace{-1.6cm}\Huge\bfseries A Linear Algebra Guide to the\\Strange Life of High-Dimensional Data\\[0.4em]
\Large Preliminary Working Report --- Questions 1, 2 \& 3}
\author{K.\ Vishnu Kumar \and Adithya Kesan Jayakanth \and Adithya Anand}
\date{27 May 2026\\[0.3em]\small\itshape Internal draft: full proofs and all details.
Questions 1, 2 (analytical) and 3 (data analysis on \texttt{mystery\_matrix1.npy});
to be compressed into the 10-page submission. Extra detail is retained here for the oral defence.}

\begin{document}
\maketitle
\thispagestyle{fancy}

% =====================================================================
\section*{Scope and status of this draft}
\addcontentsline{toc}{section}{Scope and status}

This preliminary report contains the \emph{complete} worked solutions to all three
questions: \textbf{Question 1} (Linear Algebra of Markov Matrices) and \textbf{Question 2}
(The Power Method and Singular Value Decomposition), both purely analytical with every
proof written out in full, and \textbf{Question 3} (the data analysis on
\texttt{mystery\_matrix1.npy}), now added with its full computational pipeline, figures,
and propositions. It remains a working document, not the final submission: the generic
report scaffolding (literature, sensitivity studies) is kept light, and the whole will be
compressed to the 10-page cap for submission, with the surplus detail here reserved for the
oral defence.

\textbf{Status of Question 3.} Parts (a)--(h) are complete and reproduced from
\texttt{src/task\_3.ipynb} (every figure in this section is regenerated by that notebook and
stored under \texttt{data/output/TASK 3/}). Parts (a)--(e) consolidate the earlier
\texttt{doc/task 3 anand/Q3\_report.tex}; parts (f) (Gram-matrix clustering), (g) (the
covariance-vs-Gram comparison), and (h) (the taste-tester synthesis: how many unique samples) are
written up here for the first time. Question 1's
numerical-validation parts (f) and (h), flagged as misses in the 26 May draft, are now
implemented in \texttt{src/task1/task1.ipynb} and resolved below.

Coloured tags mark every place this document departs from, or extends, the team's
handwritten notes and the raw data. Each tag has a matching entry in \texttt{changelog.md}:
\begin{itemize}[leftmargin=1.4em, itemsep=1pt]
  \item \miss{\dots} --- a part of the problem not yet attempted in the notes.
  \item \fix{\dots} --- a mathematical error in the notes or a mismatch in the prompt, corrected here.
  \item \extra{\dots} --- content added beyond the notes (an elaboration or a missing explanation).
  \item \verified{\dots} --- a numerical or algebraic claim independently checked in code.
\end{itemize}

% =====================================================================

\end{document}
